\section{Reintegrated Historical Proof Blocks (archival proof-memory)}
\label{app:formal-reintegrated}

\textit{Read this appendix as archival proof-memory and full-source preservation beneath the active main line. In earlier working states this location contained only a compact recovered fragment. To avoid archival page loss and to preserve the complete historical witness, that compact fragment is replaced here by the historical-only Part~II carried forward from the preserved v3.25.0 release line. Whenever wording in the archived text sounds stronger, flatter, or more present-tense than the normalized live line, the current status statements and routing discipline in the active main text govern.}

\begin{remark}[Preservation note]
The appendix that follows is intentionally the integrated historical witness rather than a shortened reintegrated excerpt. It preserves the older derivation layouts and source-memory material without re-importing the active current document a second time.
\end{remark}


\clearpage
\section*{Historical witness handoff}
\addcontentsline{toc}{section}{Historical witness handoff}
\begin{quote}\small
\noindent\textbf{Bundle status.} Beginning with this public split-bundle release, the active core and the integrated historical witness are carried as separate PDFs in the same release bundle. This is a readability and archival packaging change only. It does not delete the historical witness, does not change theorematic status, and does not manufacture any OP closure certificate.

\medskip
\noindent\textbf{Companion file.} The preserved historical witness is provided as\par
\texttt{Quantum\_Sphaera\_Companion\_v3\_35\_0\_historical\_witness\_release\_final.pdf}.

\medskip
\noindent\textbf{Reading rule.} The present active-core PDF is the governing current proof-line and audit interface. The historical-witness PDF is preserved proof memory, genealogy, and reconstruction support. Whenever an inherited historical block and the active core differ in theorematic status, the active core governs.
\end{quote}

\begin{thebibliography}{99}

\bibitem{brendecke2026rh}
Marc Brendecke,
\emph{On the Riemann Hypothesis: A Kernel-Theoretic Approach via Tensor Contraction and Self-Adjoint Extension},
Zenodo v4.15, March 2026.
\newblock DOI: \href{https://doi.org/10.5281/zenodo.19157158}{10.5281/zenodo.19157158}.

\bibitem{brendecke2026main}
Marc Brendecke,
\emph{Quantum Sphaera: A Cascade Framework for Quantum Gravity, the Standard Model, and the Riemann Hypothesis},
Zenodo, March 2026.
\newblock DOI: \href{https://doi.org/10.5281/zenodo.19210213}{10.5281/zenodo.19210213}.

\bibitem{brendecke2026sphaera}
Marc Brendecke,
\emph{Quantum Sphaera Companion: A Structure-First Mathematical Unfolding} (v3.35.0 pre-v4 QAM reindexing, two-sided tensor-notation, and mixed-readout release),
April 2026.
\newblock Concept DOI: \href{https://doi.org/10.5281/zenodo.19210728}{10.5281/zenodo.19210728}; version-specific DOI: \href{https://doi.org/10.5281/zenodo.19750674}{10.5281/zenodo.19750674}; prior public version DOI: \href{https://doi.org/10.5281/zenodo.19746821}{10.5281/zenodo.19746821}.

\bibitem{kim2026}
Yilwook Kim,
\emph{GUE-pacing and Energy Highway in the Riemann Lattice},
Zenodo, March 2026.
\newblock Concept DOI: \texttt{10.5281/zenodo.19112361}.

\bibitem{Riemann1859}
B.~Riemann,
\emph{\"Uber die Anzahl der Primzahlen unter einer gegebenen Gr\"osse},
Monatsberichte der Berliner Akademie, November 1859, 136--144.

\bibitem{Edwards1974}
H.~M. Edwards,
\emph{Riemann's Zeta Function},
Academic Press, New York, 1974.

\bibitem{Titchmarsh1986}
E.~C. Titchmarsh,
\emph{The Theory of the Riemann Zeta-Function},
2nd ed., revised by D.~R. Heath-Brown,
Oxford University Press, 1986.

\bibitem{Conrey2003}
J.~B. Conrey,
\emph{The Riemann Hypothesis},
Notices of the American Mathematical Society \textbf{50} (2003), no.~3, 341--353.

\bibitem{Montgomery1973}
H.~L. Montgomery,
\emph{The Pair Correlation of Zeros of the Zeta Function},
Proc. Sympos. Pure Math. \textbf{24} (1973), 181--193.

\bibitem{Odlyzko1987}
A.~M. Odlyzko,
\emph{On the Distribution of Spacings Between Zeros of the Zeta Function},
Mathematics of Computation \textbf{48} (1987), no.~177, 273--308.
\newblock DOI: \href{https://doi.org/10.1090/S0025-5718-1987-0866115-0}{10.1090/S0025-5718-1987-0866115-0}.

\bibitem{Mehta2004}
M.~L. Mehta,
\emph{Random Matrices},
3rd ed., Pure and Applied Mathematics, vol.~142,
Elsevier/Academic Press, 2004.

\bibitem{KatzSarnak1999}
N.~M. Katz and P.~Sarnak,
\emph{Zeroes of Zeta Functions and Symmetry},
Bulletin of the American Mathematical Society \textbf{36} (1999), no.~1, 1--26.
\newblock DOI: \href{https://doi.org/10.1090/S0273-0979-99-00766-1}{10.1090/S0273-0979-99-00766-1}.

\bibitem{Ahlfors2006}
L.~V. Ahlfors,
\emph{Lectures on Quasiconformal Mappings},
2nd ed., University Lecture Series, vol.~38,
American Mathematical Society, Providence, RI, 2006.
\newblock DOI: \href{https://doi.org/10.1090/ulect/038}{10.1090/ulect/038}.

\bibitem{GardinerLakic2000}
F.~P. Gardiner and N.~Lakic,
\emph{Quasiconformal Teichm\"uller Theory},
Mathematical Surveys and Monographs, vol.~76,
American Mathematical Society, Providence, RI, 2000.

\bibitem{Hubbard2006}
J.~H. Hubbard,
\emph{Teichm\"uller Theory and Applications to Geometry, Topology, and Dynamics, Volume 1: Teichm\"uller Theory},
Matrix Editions, Ithaca, NY, 2006.

\bibitem{Sullivan1985}
D.~Sullivan,
\emph{Quasiconformal Homeomorphisms and Dynamics I. Solution of the Fatou--Julia Problem on Wandering Domains},
Annals of Mathematics \textbf{122} (1985), no.~3, 401--418.
\newblock DOI: \href{https://doi.org/10.2307/1971308}{10.2307/1971308}.

\bibitem{Milnor2006}
J.~Milnor,
\emph{Dynamics in One Complex Variable},
3rd ed., Annals of Mathematics Studies, vol.~160,
Princeton University Press, 2006.

\bibitem{Hopf1931}
H.~Hopf,
\emph{\"Uber die Abbildungen der dreidimensionalen Sph\"are auf die Kugelfl\"ache},
Mathematische Annalen \textbf{104} (1931), 637--665.
\newblock DOI: \href{https://doi.org/10.1007/BF01457962}{10.1007/BF01457962}.

\bibitem{Steenrod1951}
N.~Steenrod,
\emph{The Topology of Fibre Bundles},
Princeton Mathematical Series, vol.~14,
Princeton University Press, 1951.

\bibitem{Hatcher2002}
A.~Hatcher,
\emph{Algebraic Topology},
Cambridge University Press, 2002.

\bibitem{Bryant1987}
R.~L. Bryant,
\emph{Metrics with Exceptional Holonomy},
Annals of Mathematics \textbf{126} (1987), no.~3, 525--576.
\newblock DOI: \href{https://doi.org/10.2307/1971360}{10.2307/1971360}.

\bibitem{EguchiGilkeyHanson1980}
T.~Eguchi, P.~B. Gilkey, and A.~J. Hanson,
\emph{Gravitation, Gauge Theories and Differential Geometry},
Physics Reports \textbf{66} (1980), no.~6, 213--393.
\newblock DOI: \href{https://doi.org/10.1016/0370-1573(80)90130-1}{10.1016/0370-1573(80)90130-1}.

\bibitem{Baez2002}
J.~C. Baez,
\emph{The Octonions},
Bulletin of the American Mathematical Society \textbf{39} (2002), no.~2, 145--205.
\newblock DOI: \href{https://doi.org/10.1090/S0273-0979-01-00934-X}{10.1090/S0273-0979-01-00934-X}.

\bibitem{Dixon1994}
G.~M. Dixon,
\emph{Division Algebras: Octonions, Quaternions, Complex Numbers and the Algebraic Design of Physics},
Mathematics and Its Applications, vol.~290,
Kluwer Academic Publishers, Dordrecht, 1994.
\newblock DOI: \href{https://doi.org/10.1007/978-1-4757-2315-1}{10.1007/978-1-4757-2315-1}.

\bibitem{Furey2018}
C.~Furey,
\emph{$SU(3)_C\times SU(2)_L\times U(1)_Y\left(\times U(1)_X\right)$ as a Symmetry of Division Algebraic Ladder Operators},
European Physical Journal C \textbf{78} (2018), 375.
\newblock DOI: \href{https://doi.org/10.1140/epjc/s10052-018-5844-7}{10.1140/epjc/s10052-018-5844-7}.

\bibitem{Dirac1928}
P.~A.~M. Dirac,
\emph{The Quantum Theory of the Electron},
Proceedings of the Royal Society of London A \textbf{117} (1928), no.~778, 610--624.
\newblock DOI: \href{https://doi.org/10.1098/rspa.1928.0023}{10.1098/rspa.1928.0023}.

\bibitem{YangMills1954}
C.~N. Yang and R.~L. Mills,
\emph{Conservation of Isotopic Spin and Isotopic Gauge Invariance},
Physical Review \textbf{96} (1954), 191--195.
\newblock DOI: \href{https://doi.org/10.1103/PhysRev.96.191}{10.1103/PhysRev.96.191}.

\bibitem{EnglertBrout1964}
F.~Englert and R.~Brout,
\emph{Broken Symmetry and the Mass of Gauge Vector Mesons},
Physical Review Letters \textbf{13} (1964), 321--323.
\newblock DOI: \href{https://doi.org/10.1103/PhysRevLett.13.321}{10.1103/PhysRevLett.13.321}.

\bibitem{Higgs1964}
P.~W. Higgs,
\emph{Broken Symmetries and the Masses of Gauge Bosons},
Physical Review Letters \textbf{13} (1964), 508--509.
\newblock DOI: \href{https://doi.org/10.1103/PhysRevLett.13.508}{10.1103/PhysRevLett.13.508}.

\bibitem{Cabibbo1963}
N.~Cabibbo,
\emph{Unitary Symmetry and Leptonic Decays},
Physical Review Letters \textbf{10} (1963), 531--533.
\newblock DOI: \href{https://doi.org/10.1103/PhysRevLett.10.531}{10.1103/PhysRevLett.10.531}.

\bibitem{KobayashiMaskawa1973}
M.~Kobayashi and T.~Maskawa,
\emph{CP-Violation in the Renormalizable Theory of Weak Interaction},
Progress of Theoretical Physics \textbf{49} (1973), no.~2, 652--657.
\newblock DOI: \href{https://doi.org/10.1143/PTP.49.652}{10.1143/PTP.49.652}.

\bibitem{MakiNakagawaSakata1962}
Z.~Maki, M.~Nakagawa, and S.~Sakata,
\emph{Remarks on the Unified Model of Elementary Particles},
Progress of Theoretical Physics \textbf{28} (1962), 870--880.
\newblock DOI: \href{https://doi.org/10.1143/PTP.28.870}{10.1143/PTP.28.870}.

\bibitem{Stone1930}
M.~H.~Stone,
\emph{Linear Transformations in Hilbert Space, III: Operational Methods and Group Theory},
Proceedings of the National Academy of Sciences of the USA \textbf{16} (1930), 172--175.

\bibitem{vonNeumann1931}
J.~von Neumann,
\emph{Die Eindeutigkeit der Schr\"odingerschen Operatoren},
Mathematische Annalen \textbf{104} (1931), 570--578.

\bibitem{Weil1964}
A.~Weil,
\emph{Sur certains groupes d'op\'erateurs unitaires},
Acta Mathematica \textbf{111} (1964), 143--211.
\newblock DOI: \href{https://doi.org/10.1007/BF02391012}{10.1007/BF02391012}.

\bibitem{ReedSimon1975}
M.~Reed and B.~Simon,
\emph{Methods of Modern Mathematical Physics II: Fourier Analysis, Self-Adjointness},
Academic Press, New York, 1975.

\bibitem{KadisonRingrose1997}
R.~V. Kadison and J.~R. Ringrose,
\emph{Fundamentals of the Theory of Operator Algebras, Volume I: Elementary Theory},
Graduate Studies in Mathematics, vol.~15,
American Mathematical Society, 1997.

\bibitem{Zermelo1908}
E.~Zermelo,
\emph{Untersuchungen \"uber die Grundlagen der Mengenlehre. I},
Mathematische Annalen \textbf{65} (1908), 261--281.
\newblock DOI: \href{https://doi.org/10.1007/BF01449999}{10.1007/BF01449999}.

\bibitem{Fraenkel1928}
A.~A. Fraenkel,
\emph{Einleitung in die Mengenlehre},
3rd ed., Springer, Berlin, 1928.
\newblock DOI: \href{https://doi.org/10.1007/978-3-662-42029-4}{10.1007/978-3-662-42029-4}.

\bibitem{Jech2003}
T.~Jech,
\emph{Set Theory: The Third Millennium Edition, Revised and Expanded},
Springer Monographs in Mathematics,
Springer, Berlin, 2003.

\bibitem{Kunen1980}
K.~Kunen,
\emph{Set Theory: An Introduction to Independence Proofs},
Studies in Logic and the Foundations of Mathematics, vol.~102,
North-Holland, Amsterdam, 1980.

\bibitem{Turing1936}
A.~M. Turing,
\emph{On Computable Numbers, with an Application to the Entscheidungsproblem},
Proceedings of the London Mathematical Society (2) \textbf{42} (1936), 230--265.

\bibitem{Kasiski1863}
F.~W. Kasiski,
\emph{Die Geheimschriften und die Dechiffrir-Kunst: Mit besonderer Ber\"ucksichtigung der deutschen und der franz\"osischen Sprache},
E. S. Mittler und Sohn, Berlin, 1863.

\bibitem{RivestShamirAdleman1978}
R.~L. Rivest, A.~Shamir, and L.~Adleman,
\emph{A Method for Obtaining Digital Signatures and Public-Key Cryptosystems},
Communications of the ACM \textbf{21} (1978), no.~2, 120--126.
\newblock DOI: \href{https://doi.org/10.1145/359340.359342}{10.1145/359340.359342}.

\bibitem{DiffieHellman1976}
W.~Diffie and M.~E. Hellman,
\emph{New Directions in Cryptography},
IEEE Transactions on Information Theory \textbf{22} (1976), no.~6, 644--654.
\newblock DOI: \href{https://doi.org/10.1109/TIT.1976.1055638}{10.1109/TIT.1976.1055638}.

\bibitem{Norris1997}
J.~R. Norris,
\emph{Markov Chains},
Cambridge Series in Statistical and Probabilistic Mathematics, vol.~2,
Cambridge University Press, 1997.

\bibitem{Glasserman2004}
P.~Glasserman,
\emph{Monte Carlo Methods in Financial Engineering},
Stochastic Modelling and Applied Probability, vol.~53,
Springer, New York, 2004.
\newblock DOI: \href{https://doi.org/10.1007/978-0-387-21617-1}{10.1007/978-0-387-21617-1}.

\bibitem{Bollerslev1986}
T.~Bollerslev,
\emph{Generalized Autoregressive Conditional Heteroskedasticity},
Journal of Econometrics \textbf{31} (1986), no.~3, 307--327.
\newblock DOI: \href{https://doi.org/10.1016/0304-4076(86)90063-1}{10.1016/0304-4076(86)90063-1}.

\bibitem{Knuth1997}
D.~E. Knuth,
\emph{The Art of Computer Programming, Vol.~1: Fundamental Algorithms},
3rd ed., Addison--Wesley, 1997.

\bibitem{Rojas1998}
Ra\'ul Rojas,
\emph{How to Make Zuse's Z3 a Universal Computer},
IEEE Annals of the History of Computing \textbf{20} (1998), no.~3, 51--54.

\bibitem{Copeland2006}
B.~J. Copeland (ed.),
\emph{Colossus: The Secrets of Bletchley Park's Codebreaking Computers},
Oxford University Press, 2006.

\bibitem{Zuse2016}
H.~Zuse,
\emph{Konrad Zuse's Computer Z3},
historical overview, 2016.

\bibitem{MenezesOorschotVanstone1996}
A.~J. Menezes, P.~C. van Oorschot, and S.~A. Vanstone,
\emph{Handbook of Applied Cryptography},
CRC Press, 1996.

\bibitem{GoldwasserBellare2008}
S.~Goldwasser and M.~Bellare,
\emph{Lecture Notes on Cryptography},
MIT/UCSD lecture notes, 2008.

\bibitem{PeskinSchroeder1995}
M.~E. Peskin and D.~V. Schroeder,
\emph{An Introduction to Quantum Field Theory},
Westview Press, 1995.

\bibitem{Weinberg1995}
S.~Weinberg,
\emph{The Quantum Theory of Fields, Volume I: Foundations},
Cambridge University Press, 1995.

\bibitem{FultonHarris1991}
W.~Fulton and J.~Harris,
\emph{Representation Theory: A First Course},
Graduate Texts in Mathematics, vol.~129,
Springer, 1991.

\bibitem{Nakahara2003}
M.~Nakahara,
\emph{Geometry, Topology and Physics},
2nd ed., Institute of Physics Publishing, 2003.

\bibitem{Besse1987}
A.~L. Besse,
\emph{Einstein Manifolds},
Ergebnisse der Mathematik und ihrer Grenzgebiete, vol.~10,
Springer, 1987.

\bibitem{Joyce2000}
D.~D. Joyce,
\emph{Compact Manifolds with Special Holonomy},
Oxford University Press, 2000.

\end{thebibliography}
